
% !TeX root=../main.tex
\chapter{مروری بر مطالعات انجام شده}
%\thispagestyle{empty} 
\section{مقدمه و تعریف مسئله ناوبری خودران اسکوتر}

ناوبری خودران ربات‌های متحرک چرخ‌دار در محیط‌های واقعی، به‌ویژه محیط‌های دینامیکی، یکی از مسائل بنیادی و در عین حال چالش‌برانگیز در حوزه رباتیک به شمار می‌رود. با گسترش کاربرد ربات‌های خدمت‌رسان، وسایل نقلیه سبک خودران و سامانه‌های حمل‌ونقل هوشمند، نیاز به روش‌هایی که بتوانند حرکت ایمن، روان و قابل‌اعتماد را بدون دخالت مستقیم انسان تضمین کنند، بیش از پیش احساس می‌شود. در این میان، ناوبری در محیط‌های دینامیکی از اهمیت ویژه‌ای برخوردار است، زیرا ربات باید در حضور تغییرات مداوم محیط و رفتارهای غیرقابل پیش‌بینی تصمیم‌گیری کند%
\cite{Zhu2021}.

محیط دینامیکی در مسئله ناوبری به محیطی اطلاق می‌شود که در آن موانع ثابت تنها بخش کوچکی از چالش را تشکیل می‌دهند و بخش اصلی پیچیدگی ناشی از حضور موانع متحرک با سرعت و مسیر متغیر، عامل‌های تصمیم‌گیر مستقل (مانند انسان‌ها یا سایر ربات‌ها) و تغییرات ساختاری موقتی محیط است. در چنین شرایطی، ربات نه‌تنها باید از برخورد اجتناب کند، بلکه باید بتواند پیامدهای زمانی تصمیم‌های خود را نیز در نظر بگیرد؛ به‌عنوان مثال، پیش‌بینی کند که یک مانع متحرک در آینده نزدیک در چه موقعیتی قرار خواهد گرفت و آیا مسیر انتخاب‌شده همچنان ایمن باقی می‌ماند یا خیر. این مسئله، فرآیند تصمیم‌گیری را از یک مسئله صرفاً هندسی به یک مسئله پویا و وابسته به زمان تبدیل می‌کند%
\cite{s25113394}.

روش‌های کلاسیک ناوبری که عمدتاً بر پایه مدل‌های صریح، قوانین از پیش تعریف‌شده و توابع هزینه ثابت طراحی شده‌اند، در مواجهه با چنین شرایطی با محدودیت‌های جدی روبه‌رو می‌شوند. بسیاری از این روش‌ها فرض می‌کنند که محیط یا کاملاً ایستا است یا تغییرات آن به‌صورت ساده و قابل پیش‌بینی قابل مدل‌سازی است. در نتیجه، در حضور موانع با رفتار غیرقابل پیش‌بینی یا تعاملات چندعاملی، نیاز به تنظیم دستی پارامترها افزایش یافته و عملکرد الگوریتم به‌شدت به شرایط خاص محیط وابسته می‌شود. این وابستگی، تعمیم‌پذیری روش‌های کلاسیک را در محیط‌های واقعی و متغیر محدود می‌کند%
\cite{s23249672}.

علاوه بر دینامیکی بودن محیط، ویژگی‌های فیزیکی و دینامیکی پلتفرم نیز نقش تعیین‌کننده‌ای در پیچیدگی مسئله ناوبری دارند. اسکوتر چهار چرخ محرک، در مقایسه با ربات‌های دیفرانسیلی سبک آزمایشگاهی، دارای جرم بالاتر، اینرسی بیشتر، فاصله ترمز طولانی‌تر و محدودیت‌های حرکتی متفاوتی است. این ویژگی‌ها باعث می‌شوند که تصمیم‌گیری‌های ناوبری حساس‌تر بوده و خطاهای کوچک در انتخاب مسیر یا سرعت، پیامدهای جدی‌تری از نظر ایمنی و پایداری حرکت به‌دنبال داشته باشند. از این رو، بسیاری از فرضیات ساده‌کننده‌ای که در ناوبری ربات‌های کوچک به‌کار می‌روند، در مورد چنین پلتفرم‌هایی معتبر نخواهند بود.

در سال‌های اخیر، پیشرفت‌های قابل توجه در «یادگیری تقویتی عمیق%
\LTRfootnote{Deep Reinforcement Learning -- DRL}»
افق‌های جدیدی را برای حل مسئله ناوبری در محیط‌های پیچیده و دینامیکی گشوده است. برخلاف روش‌های کلاسیک، DRL این امکان را فراهم می‌کند که ربات از طریق تعامل مستقیم با محیط و بدون نیاز به مدل‌سازی صریح تمامی حالات، سیاستی را بیاموزد که متناسب با شرایط واقعی و متغیر محیط تصمیم‌گیری کند. مرورهای جامع منتشرشده در سال‌های اخیر نشان می‌دهند که DRL به یکی از جریان‌های غالب در حل مسئله ناوبری ربات‌های متحرک در محیط‌های دینامیکی تبدیل شده است%
\cite{Zhu2021,s25113394}.

با وجود این پیشرفت‌ها، بخش قابل توجهی از مطالعات موجود همچنان بر روی پلتفرم‌های ساده، محیط‌های محدود یا سناریوهای با پیچیدگی کاهش‌یافته متمرکز هستند و بررسی ناوبری مبتنی بر DRL برای پلتفرم‌های واقعی‌تر با محدودیت‌های دینامیکی جدی‌تر، به‌ویژه در محیط‌های دینامیکی و شلوغ، کمتر مورد توجه قرار گرفته است. این خلأ پژوهشی، ضرورت توسعه و ارزیابی روش‌های یادگیری تقویتی عمیق برای ناوبری خودران اسکوتر چهار چرخ محرک در محیط‌های دینامیکی را برجسته می‌سازد.

بر این اساس، هدف این پایان‌نامه طراحی و بررسی یک سامانه ناوبری خودران مبتنی بر یادگیری تقویتی عمیق است که بتواند با در نظر گرفتن ویژگی‌های دینامیکی اسکوتر و ماهیت متغیر محیط، تصمیم‌های حرکتی ایمن، هموار و کارا را در حضور موانع متحرک و شرایط غیرقابل پیش‌بینی اتخاذ نماید.

\section{روش‌های کلاسیک ناوبری ربات‌های متحرک}

پیش از گسترش رویکردهای مبتنی بر یادگیری، مسئله ناوبری ربات‌های متحرک عمدتاً با استفاده از روش‌های کلاسیک حل می‌شد. این روش‌ها معمولاً بر پایه مدل‌های صریح محیط، قوانین هندسی، جست‌وجو در فضای حالت، توابع هزینه یا حل مسائل بهینه‌سازی طراحی می‌شوند. مزیت اصلی آن‌ها شفافیت، قابلیت تحلیل، امکان پیاده‌سازی بلادرنگ و رفتار نسبتاً قابل پیش‌بینی است. به همین دلیل، بسیاری از سامانه‌های عملی ناوبری هنوز نیز از ترکیبی از این روش‌ها، به‌ویژه برای برنامه‌ریزی مسیر کلی و اجتناب محلی از برخورد، استفاده می‌کنند.

با این حال، روش‌های کلاسیک معمولاً بر فرض‌هایی مانند دسترسی به نقشه نسبتاً دقیق، قابلیت پیش‌بینی محیط، مدل دینامیکی مشخص و تنظیم مناسب پارامترها متکی هستند. این فرض‌ها در محیط‌های دینامیکی، شلوغ و ناقصاً مشاهده‌پذیر همواره برقرار نیستند. از این رو، بررسی این روش‌ها برای روشن شدن محدودیت‌های آن‌ها و توجیه حرکت به سمت رویکردهای یادگیری‌محور ضروری است.

\subsection{برنامه‌ریزی مسیر سراسری مبتنی بر گراف و جست‌وجو}

یکی از قدیمی‌ترین و پرکاربردترین خانواده‌های روش‌های کلاسیک ناوبری، روش‌های مبتنی بر گراف و جست‌وجو هستند. در این رویکردها، محیط معمولاً به‌صورت یک شبکه، گراف یا نقشه اشغال نمایش داده می‌شود و مسئله ناوبری به یافتن مسیر کم‌هزینه میان موقعیت اولیه و هدف تبدیل می‌گردد. الگوریتم‌هایی مانند \lr{Dijkstra}، \lr{A*}، \lr{D*} و \lr{D* Lite} از شناخته‌شده‌ترین روش‌های این دسته هستند%
\cite{Dijkstra1959,Hart1968AStar,koenig2005dstar}.
در این روش‌ها، هزینه مسیر می‌تواند بر اساس طول مسیر، فاصله از موانع، زمان حرکت یا ترکیبی از معیارهای هندسی تعریف شود.

روش‌های جست‌وجومحور به‌ویژه برای محیط‌های ایستا و دارای نقشه مشخص مناسب هستند، زیرا مسئله ناوبری را به جست‌وجوی مسیر کم‌هزینه در یک نمایش گرافی از محیط تبدیل می‌کنند. برای مثال، الگوریتم \lr{A*} با استفاده از تابع ابتکاری می‌تواند فرآیند جست‌وجو را نسبت به روش‌های جست‌وجوی کور هدفمندتر کند. همچنین روش‌هایی مانند \lr{D*} و \lr{D* Lite} برای بازبرنامه‌ریزی مسیر در محیط‌هایی که بخشی از نقشه به‌تدریج آشکار می‌شود، توسعه یافته‌اند و در ربات‌های متحرک کاربرد زیادی دارند%
\cite{koenig2005dstar,ferguson2008motion_planning}.

با وجود این مزایا، محدودیت اصلی این دسته از روش‌ها آن است که معمولاً مسیر را در سطح هندسی و سراسری تولید می‌کنند و به‌تنهایی تضمین نمی‌کنند که مسیر تولیدشده با دینامیک واقعی پلتفرم، محدودیت شتاب، محدودیت ترمز و قیود ایمنی لحظه‌ای سازگار باشد. علاوه بر این، در محیط‌های دینامیکی که موانع متحرک به‌صورت پیوسته موقعیت خود را تغییر می‌دهند، مسیر سراسری ممکن است به‌سرعت نامعتبر شود و سیستم به بازبرنامه‌ریزی مکرر یا ترکیب با یک برنامه‌ریز محلی نیاز پیدا کند%
\cite{LaValle2006Planning,choset2005principles}.
این مسئله برای پلتفرم‌هایی مانند اسکوتر چهار چرخ محرک که دارای اینرسی و فاصله توقف قابل توجه هستند، اهمیت بیشتری دارد.

\subsection{روش‌های نمونه‌برداری در فضای پیکربندی}

دسته دیگری از روش‌های کلاسیک ناوبری، روش‌های نمونه‌برداری در فضای پیکربندی هستند. در این رویکردها، به‌جای جست‌وجوی کامل یا گسسته‌سازی منظم فضای پیکربندی، مجموعه‌ای از پیکربندی‌های مجاز نمونه‌برداری شده و سپس به‌صورت گراف یا درخت به یکدیگر متصل می‌شوند. روش‌هایی مانند \lr{PRM}، \lr{RRT} و \lr{RRT*} از نمونه‌های شناخته‌شده این خانواده هستند%
\cite{Kavraki1996PRM,LaValle1998RRT,Karaman2011RRTstar}.
این الگوریتم‌ها به‌ویژه در فضاهای پیکربندی با ابعاد بالا یا محیط‌هایی که تولید شبکه منظم در آن‌ها دشوار و پرهزینه است، کاربرد زیادی دارند%
\cite{LaValle2006Planning,Choset2005Principles}.

مزیت اصلی روش‌های نمونه‌برداری، توانایی آن‌ها در یافتن مسیر در محیط‌های پیچیده بدون نیاز به گسسته‌سازی کامل فضای پیکربندی است. برای نمونه، روش \lr{RRT} با رشد سریع در فضای آزاد، امکان یافتن مسیرهای قابل اجرا را در زمان نسبتاً کوتاه فراهم می‌کند و نسخه‌هایی مانند \lr{RRT*} تلاش می‌کنند کیفیت مسیر را از نظر بهینگی بهبود دهند%
\cite{LaValle1998RRT,Karaman2011RRTstar}.
با این حال، مسیرهای تولیدشده توسط این روش‌ها معمولاً خام، ناصاف و نیازمند هموارسازی یا پردازش تکمیلی هستند. همچنین، اگرچه این روش‌ها برای برنامه‌ریزی مسیر مفیدند، اما در حالت پایه برای محیط‌های به‌شدت دینامیکی و تصمیم‌گیری بلادرنگ در حضور موانع متحرک کافی نیستند%
\cite{LaValle2006Planning,Choset2005Principles}.

از منظر پژوهش حاضر، روش‌های نمونه‌برداری می‌توانند برای تولید مسیر اولیه یا مسیر مرجع مفید باشند، اما به‌تنهایی مسئله ناوبری پیوسته و ایمن اسکوتر در محیط دینامیکی را حل نمی‌کنند. دلیل اصلی آن است که اسکوتر نه‌تنها باید از نظر هندسی به هدف برسد، بلکه باید فرمان‌هایی تولید کند که با محدودیت‌های حرکتی، شتاب، ترمز، پایداری و نرمی حرکت نیز سازگار باشند.

\subsection{میدان‌های پتانسیل مصنوعی و مسئله کمینه‌های محلی}

یکی از نخستین رویکردهای مطرح در ناوبری ربات‌های متحرک، روش «میدان‌های پتانسیل مصنوعی%
\LTRfootnote{\lr{Artificial Potential Fields -- APF}}»
است. در این روش، هدف به‌عنوان یک پتانسیل جاذب و موانع به‌عنوان پتانسیل‌های دافع مدل می‌شوند و ربات با حرکت در راستای کاهش تابع پتانسیل به سمت هدف هدایت می‌شود. سادگی محاسبات، پیاده‌سازی آسان و امکان اجرای بلادرنگ از مهم‌ترین مزایای این روش به‌شمار می‌روند%
\cite{khatib1986}.

با وجود این مزایا، یکی از مشکلات بنیادین میدان‌های پتانسیل، گیر افتادن ربات در نقاط «کمینه محلی%
\LTRfootnote{\lr{Local Minima}}»
است. در چنین نقاطی، اثر پتانسیل جاذب هدف و پتانسیل‌های دافع موانع می‌تواند به تعادل برسد و گرادیان مؤثر میدان پتانسیل صفر یا بسیار کوچک شود؛ در نتیجه، ربات ممکن است بدون رسیدن به هدف متوقف شود. این پدیده به‌ویژه در محیط‌های دارای موانع مقعر، راهروهای باریک یا چینش‌های پیچیده موانع تشدید می‌شود. افزون بر این، روش میدان پتانسیل در حالت پایه معمولاً قیود دینامیکی پلتفرم، محدودیت ترمز، محدودیت شتاب و رفتار زمانی موانع متحرک را به‌صورت صریح در نظر نمی‌گیرد%
\cite{Koren1991PotentialFields,LaValle2006Planning,Choset2005Principles}.
بنابراین، اگرچه این روش از نظر مفهومی ساده و قابل فهم است، اما برای ناوبری ایمن در محیط‌های دینامیکی و برای پلتفرم‌هایی با دینامیک جدی، به‌تنهایی کافی نیست.

\subsection{روش‌های اجتناب از برخورد مبتنی بر سرعت}

برای مواجهه با موانع متحرک، خانواده‌ای از روش‌ها توسعه یافته‌اند که به‌جای تمرکز صرف بر موقعیت هندسی، فضای سرعت‌های مجاز و غیرمجاز را تحلیل می‌کنند. روش‌هایی مانند «مانع سرعت%
\LTRfootnote{\lr{Velocity Obstacle -- VO}}»،
«مانع سرعت متقابل%
\LTRfootnote{\lr{Reciprocal Velocity Obstacle -- RVO}}»
و «اجتناب بهینه متقابل از برخورد%
\LTRfootnote{\lr{Optimal Reciprocal Collision Avoidance -- ORCA}}»
در این دسته قرار می‌گیرند%
\cite{Fiorini1998VO,VanDenBerg2008RVO,VanDenBerg2011ORCA}.
ایده اصلی این روش‌ها آن است که با توجه به سرعت نسبی ربات و موانع متحرک، می‌توان نواحی‌ای از فضای سرعت را مشخص کرد که انتخاب آن‌ها در یک افق زمانی محدود به برخورد منجر می‌شود. سپس فرمان حرکتی ربات به‌گونه‌ای انتخاب می‌شود که از این نواحی خطرناک در فضای سرعت دوری کند.

این دسته از روش‌ها برای محیط‌های چندعاملی و موانع متحرک اهمیت زیادی دارند، زیرا مسئله برخورد را به‌صورت زمانی بررسی می‌کنند و فقط به فاصله لحظه‌ای از مانع وابسته نیستند. با این حال، عملکرد آن‌ها معمولاً به فرض‌هایی مانند حرکت خطی کوتاه‌مدت عامل‌ها، واکنش متقابل سایر عامل‌ها و دقت در تخمین سرعت موانع وابسته است. در محیط‌های واقعی، به‌ویژه در حضور انسان‌ها، رفتار موانع لزوماً خطی، منطقی یا قابل پیش‌بینی نیست. همچنین، این روش‌ها معمولاً مسئله را در سطح اجتناب محلی از برخورد حل می‌کنند و برای تصمیم‌گیری بلندمدت، انتخاب مسیر کلی و لحاظ کامل قیود دینامیکی پلتفرم نیازمند ترکیب با سایر روش‌ها هستند%
\cite{Fiorini1998VO,VanDenBerg2008RVO,VanDenBerg2011ORCA,LaValle2006Planning}.

\subsection{روش پنجره دینامیکی و برنامه‌ریزی در فضای سرعت}

روش «پنجره دینامیکی%
\LTRfootnote{\lr{Dynamic Window Approach -- DWA}}»
یکی از شناخته‌شده‌ترین روش‌های برنامه‌ریزی محلی برای ربات‌های متحرک است. برخلاف روش‌های مبتنی بر موقعیت، \lr{DWA} مستقیماً در فضای سرعت خطی و زاویه‌ای ربات جست‌وجو می‌کند و تنها سرعت‌هایی را در نظر می‌گیرد که با قیود دینامیکی پلتفرم، از جمله حداکثر سرعت، شتاب و توان توقف، سازگار باشند. برای هر زوج سرعت، یک مسیر کوتاه‌مدت شبیه‌سازی شده و سپس با استفاده از تابع هزینه، مناسب‌ترین فرمان انتخاب می‌شود%
\cite{fox1997dwa}.

محبوبیت \lr{DWA} ناشی از امکان اجرای بلادرنگ، سادگی نسبی و قابلیت لحاظ برخی قیود حرکتی است. با این حال، عملکرد این روش به‌شدت به تنظیم ضرایب تابع هزینه وابسته است؛ ضرایبی که معمولاً وزن‌هایی برای حرکت به سمت هدف، حفظ فاصله ایمن از موانع و انتخاب سرعت مناسب تعیین می‌کنند. در برخی پیاده‌سازی‌ها، معیارهایی مانند همواری حرکت یا کاهش تغییرات ناگهانی فرمان نیز به تابع هزینه افزوده می‌شود. در صورت تغییر محیط، سرعت ربات یا ویژگی‌های دینامیکی پلتفرم، این ضرایب ممکن است نیازمند تنظیم مجدد باشند%
\cite{fox1997dwa,LaValle2006Planning,Choset2005Principles}.

افزون بر این، افق تصمیم‌گیری \lr{DWA} کوتاه‌مدت است؛ بنابراین این روش معمولاً تصمیم محلی مناسبی تولید می‌کند، اما به‌تنهایی تضمین‌کننده تصمیم‌گیری بلندمدت یا انتخاب مسیر کلی بهینه نیست. به همین دلیل، در محیط‌های شلوغ و دینامیکی ممکن است رفتارهایی مانند توقف‌های مکرر، حرکت محافظه‌کارانه یا نوسان میان چند فرمان نزدیک به هم رخ دهد. این مسئله برای اسکوتر چهار چرخ محرک مهم‌تر است، زیرا تغییرات ناگهانی فرمان یا توقف‌های مکرر می‌تواند موجب کاهش نرمی حرکت، افت راحتی سرنشین و افزایش ریسک ایمنی شود.

\subsection{برنامه‌ریزی مسیر محلی مبتنی بر بهینه‌سازی و روش \lr{TEB}}

در سامانه‌های رباتیک عملی، علاوه بر روش‌هایی مانند \lr{DWA}، از برنامه‌ریزهای محلی مبتنی بر بهینه‌سازی نیز استفاده می‌شود. یکی از نمونه‌های رایج این دسته، روش «نوار کشسان زمان‌بندی‌شده%
\LTRfootnote{\lr{Timed Elastic Band -- TEB}}»
است. این روش را می‌توان توسعه‌ای از ایده نوارهای کشسان در برنامه‌ریزی حرکت دانست که در آن مسیر محلی به‌صورت مجموعه‌ای از حالت‌های متوالی همراه با اطلاعات زمانی نمایش داده می‌شود%
\cite{Quinlan1993ElasticBands,Rosmann2013TEB,Rosmann2015TEB}.
سپس با در نظر گرفتن معیارهایی مانند فاصله از موانع، محدودیت سرعت، محدودیت شتاب، زمان حرکت و برخی قیود سینماتیکی، یک مسئله بهینه‌سازی محلی حل می‌گردد. مزیت اصلی این روش آن است که برخلاف مسیرهای صرفاً هندسی، زمان‌بندی حرکت و برخی محدودیت‌های حرکتی را نیز در فرآیند برنامه‌ریزی وارد می‌کند.

با وجود این مزایا، روش‌هایی مانند \lr{TEB} نیز معمولاً به تابع هزینه، وزن‌دهی مناسب قیود، کیفیت اطلاعات محیط و سرعت به‌روزرسانی نقشه محلی وابسته‌اند. در محیط‌های دینامیکی، عملکرد این روش به توانایی سامانه در تخمین یا پیش‌بینی حرکت موانع و بازبهینه‌سازی سریع مسیر بستگی دارد. اگر رفتار موانع غیرقابل پیش‌بینی باشد یا داده‌های حسگری دارای تأخیر و نویز باشند، کیفیت مسیر بهینه‌شده و قابلیت اطمینان تصمیم محلی کاهش می‌یابد. بنابراین، اگرچه \lr{TEB} نسبت به برنامه‌ریزی هندسی ساده پیشرفته‌تر است و برای برنامه‌ریزی محلی ربات‌های متحرک کاربرد زیادی دارد، اما در محیط‌های بسیار پویا، شلوغ و ناقصاً مشاهده‌پذیر همچنان با محدودیت مواجه می‌شود%
\cite{Rosmann2013TEB,Rosmann2015TEB,LaValle2006Planning}.

\subsection{کنترل پیش‌بینانه مبتنی بر مدل و برنامه‌ریزی بهینه}

«کنترل پیش‌بینانه مبتنی بر مدل%
\LTRfootnote{\lr{Model Predictive Control -- MPC}}»
به‌عنوان یک چارچوب قدرتمند برای لحاظ کردن قیود حرکتی، قیود کنترلی و اهداف بهینه‌سازی در سامانه‌های دینامیکی مطرح شده است. در روش \lr{MPC}، یک مدل از دینامیک سیستم در افق زمانی محدود در نظر گرفته می‌شود و با حل یک مسئله بهینه‌سازی مقید، توالی مناسبی از فرمان‌های کنترلی تولید می‌گردد. سپس تنها فرمان اول اعمال شده و در گام زمانی بعد، مسئله با استفاده از اطلاعات جدید دوباره حل می‌شود. این سازوکار بازخوردی باعث می‌شود \lr{MPC} بتواند در صورت دسترسی به مدل و اطلاعات به‌روز، نسبت به تغییرات وضعیت سیستم و محیط واکنش نشان دهد%
\cite{rawlings2009mpc,Bemporad1999}.

مزیت اصلی \lr{MPC} توانایی آن در لحاظ هم‌زمان قیودی مانند محدودیت سرعت، محدودیت شتاب، محدودیت ورودی‌های کنترلی، فاصله ایمن از موانع و اهداف حرکتی است. این ویژگی برای پلتفرم‌هایی مانند اسکوتر چهار چرخ محرک اهمیت زیادی دارد، زیرا این نوع پلتفرم نسبت به ربات‌های کوچک آزمایشگاهی دارای اینرسی بیشتر، محدودیت توقف جدی‌تر و حساسیت بالاتر نسبت به تغییرات ناگهانی فرمان است. با این حال، عملکرد \lr{MPC} به دقت مدل دینامیکی، کیفیت پیش‌بینی وضعیت محیط و تنظیم مناسب تابع هزینه وابسته است. علاوه بر این، حل مسئله بهینه‌سازی در هر گام زمانی می‌تواند هزینه محاسباتی قابل توجهی ایجاد کند، به‌ویژه زمانی که قیود غیرخطی، موانع متحرک یا افق پیش‌بینی بلندتر در نظر گرفته شوند.

نسخه‌های مقاوم‌تر \lr{MPC}، مانند \lr{Tube MPC}، می‌توانند تا حدی اثر عدم قطعیت‌ها و اغتشاش‌های محدود را در نظر بگیرند. با این حال، طراحی قیود مقاوم، انتخاب ناحیه‌های مجاز، تنظیم وزن‌های تابع هزینه و تضمین عملکرد مناسب در سناریوهای متنوع همچنان فرآیندی پیچیده و وابسته به مدل و دانش طراح باقی می‌ماند%
\cite{Mayne2005TubeMPC}.
از این رو، هرچند \lr{MPC} از نظر نظری و عملی یکی از قوی‌ترین روش‌های کلاسیک برای کنترل و برنامه‌ریزی مقید محسوب می‌شود، اما استفاده از آن به‌عنوان راهکار مستقل برای ناوبری کاملاً تطبیقی در محیط‌های دینامیکی، شلوغ و ناقصاً قابل پیش‌بینی با چالش‌های مهمی همراه است.

\subsection{جمع‌بندی محدودیت‌های روش‌های کلاسیک}

مرور روش‌های کلاسیک نشان می‌دهد که هر خانواده از الگوریتم‌ها بخشی از مسئله ناوبری را به‌خوبی پوشش می‌دهد، اما هیچ‌کدام به‌تنهایی پاسخ کامل و عمومی برای ناوبری در محیط‌های دینامیکی و پیچیده ارائه نمی‌کنند. روش‌های جست‌وجومحور و نمونه‌برداری برای تولید مسیر سراسری مفیدند، اما معمولاً دینامیک پلتفرم و رفتار زمانی موانع را به‌صورت کامل لحاظ نمی‌کنند%
\cite{s25113394,alvarez2024DRLsurvey}. %
 روش‌های میدان پتانسیل ساده و سریع هستند، اما با مشکل کمینه محلی و رفتارهای ناپایدار مواجه می‌شوند. روش‌های مبتنی بر سرعت مانند \lr{DWA}، \lr{Velocity Obstacle} و \lr{ORCA} برای اجتناب محلی از برخورد مناسب‌اند، اما افق تصمیم‌گیری محدود و وابستگی به فرضیات ساده درباره حرکت موانع دارند. روش‌های بهینه‌سازی‌محور مانند \lr{TEB} و \lr{MPC} نیز امکان لحاظ قیود حرکتی را فراهم می‌کنند، اما به مدل، تابع هزینه و تنظیم دقیق پارامترها وابسته‌اند.

این محدودیت‌ها در مسئله ناوبری اسکوتر چهار چرخ محرک برجسته‌تر می‌شوند، زیرا تصمیم‌های حرکتی باید علاوه بر اجتناب از برخورد، با قیود دینامیکی، اینرسی، محدودیت شتاب و ترمز، نرمی حرکت و ایمنی سرنشین نیز سازگار باشند. بنابراین، مسئله پژوهش حاضر صرفاً یافتن یک مسیر هندسی از مبدأ به هدف نیست، بلکه یادگیری یا تولید سیاستی است که بتواند در محیط دینامیکی، فرمان‌های پیوسته، ایمن و قابل اجرا تولید کند.

\section{روش‌های یادگیری ماشین در ناوبری ربات‌های متحرک}

در مسائل ناوبری ربات‌های متحرک، هدف اصلی تولید تصمیم‌های حرکتی مناسب بر اساس اطلاعات حسگری و وضعیت محیط است. در رویکردهای کلاسیک، این تصمیم‌ها معمولاً بر پایه مدل‌های تحلیلی، قوانین از پیش تعریف‌شده و توابع هزینه طراحی می‌شوند. با این حال، در بسیاری از سناریوهای واقعی، مدل‌سازی دقیق محیط و پیش‌بینی رفتار تمامی موانع، به‌ویژه در محیط‌های دینامیکی و شلوغ، بسیار دشوار یا حتی غیرممکن است. در چنین شرایطی، روش‌های مبتنی بر یادگیری ماشین به‌عنوان ابزاری برای استخراج الگوهای تصمیم‌گیری از داده و تجربه مطرح می‌شوند.

یادگیری ماشین را می‌توان به‌طور کلی به‌عنوان مجموعه‌ای از روش‌ها تعریف کرد که در آن‌ها سیستم، به‌جای اتکا به قوانین صریح یا مدل‌های از پیش تعیین‌شده، رابطه میان ورودی‌ها و خروجی‌ها را مستقیماً از داده‌ها می‌آموزد. در مسئله ناوبری، این ورودی‌ها می‌توانند شامل داده‌های حسگرهایی مانند لایدار، دوربین یا سنسورهای اینرسی باشند و خروجی‌ها معمولاً به‌صورت فرمان‌های حرکتی یا مسیرهای پیشنهادی تعریف می‌شوند. بدین ترتیب، یادگیری ماشین این امکان را فراهم می‌کند که ربات رفتار مناسب را بر اساس تجربه‌های قبلی و مشاهده شرایط مختلف محیطی استخراج کند%
\cite{Zhu2021}. %

یکی از انگیزه‌های اصلی استفاده از یادگیری ماشین در ناوبری، کاهش وابستگی به مدل‌سازی دقیق و تنظیم دستی پارامترهاست. در روش‌های کلاسیک، عملکرد الگوریتم‌ها اغلب به انتخاب ضرایب توابع هزینه یا قوانین کنترلی وابسته است و تغییر شرایط محیطی می‌تواند نیازمند تنظیم مجدد این پارامترها باشد. در مقابل، روش‌های مبتنی بر یادگیری ماشین قادرند این وابستگی را تا حدی کاهش داده و رفتار مناسب را مستقیماً از داده‌های حاصل از تعامل با محیط استخراج کنند. این ویژگی به‌ویژه در محیط‌هایی با تغییرات مداوم، عدم قطعیت بالا و رفتارهای غیرقابل پیش‌بینی اهمیت پیدا می‌کند.

با وجود این مزایا، باید توجه داشت که یادگیری ماشین یک چارچوب واحد و یکپارچه نیست، بلکه شامل خانواده‌های مختلفی از روش‌ها با فرضیات و قابلیت‌های متفاوت است. برخی از این روش‌ها بیشتر برای پیش‌بینی یا استخراج الگو از داده مناسب‌اند، در حالی که برخی دیگر امکان تصمیم‌گیری و بهینه‌سازی رفتار را در طول زمان فراهم می‌کنند. از این رو، در کاربرد ناوبری ربات‌های متحرک، انتخاب نوع روش یادگیری ماشین نقش تعیین‌کننده‌ای در میزان موفقیت سیستم دارد. در ادامه، انواع اصلی روش‌های مبتنی بر یادگیری ماشین که در حل مسائل ناوبری مورد استفاده قرار گرفته‌اند بررسی می‌شوند و مزایا و محدودیت‌های هر یک از آن‌ها در چارچوب ناوبری ربات‌های متحرک مورد بحث قرار می‌گیرد.

\subsection{یادگیری نظارتی در ناوبری ربات‌های متحرک}

«یادگیری نظارتی%
\LTRfootnote{Supervised Learning}»
یکی از شاخه‌های اصلی یادگیری ماشین است که در آن، مدل با استفاده از داده‌های برچسب‌خورده، نگاشتی میان ورودی‌ها و خروجی‌ها را می‌آموزد. در این چارچوب، خروجی مطلوب برای هر ورودی از پیش مشخص شده و هدف یادگیری، بازتولید این رفتار مرجع با کمینه‌کردن خطای پیش‌بینی است. در مسئله ناوبری ربات‌های متحرک، ورودی‌ها معمولاً شامل داده‌های حسگری مانند لایدار یا تصویر دوربین و اطلاعات هدف هستند و خروجی‌ها می‌توانند به‌صورت مسیر محلی یا فرمان‌های حرکتی بیان شوند.

در سال‌های اخیر، یادگیری نظارتی به‌ویژه در قالب رویکردهای انتها به انتها برای ناوبری ربات‌ها مورد توجه قرار گرفته است. در این روش‌ها، فرآیند تصمیم‌گیری ناوبری به‌صورت یک نگاشت مستقیم از داده‌های ادراکی به فرمان‌های کنترلی فرموله می‌شود. برای نمونه، «پفایفر%
\LTRfootnote{Pfeiffer}»%
 و همکاران یک چارچوب داده‌محور ارائه کرده‌اند که در آن داده‌های لایدار و موقعیت هدف به‌عنوان ورودی شبکه عصبی در نظر گرفته شده و خروجی شبکه، سرعت خطی و زاویه‌ای ربات را تعیین می‌کند. داده‌های آموزشی این مدل از اجرای یک سیاست مرجع مبتنی بر برنامه‌ریزی کلاسیک جمع‌آوری شده و شبکه عصبی رفتار این سیاست را تقلید می‌کند%
\cite{pfeiffer2018nav}.

مزیت اصلی چنین رویکردهایی، سادگی چارچوب تصمیم‌گیری و سرعت بالای اجرا پس از آموزش است. از آنجا که تصمیم‌گیری به‌صورت مستقیم انجام می‌شود و نیاز به حل مسائل بهینه‌سازی برخط وجود ندارد، این روش‌ها برای کاربردهای بلادرنگ جذاب هستند. همچنین، یادگیری نظارتی امکان استفاده از داده‌های انسانی یا الگوریتم‌های مرجع را فراهم می‌کند و می‌تواند به تولید مسیرهای هموار و قابل‌قبول در محیط‌های نسبتاً ساخت‌یافته منجر شود.

با وجود این مزایا، وابستگی شدید یادگیری نظارتی به کیفیت و پوشش داده‌های آموزشی، یکی از محدودیت‌های اساسی آن در ناوبری ربات‌های متحرک محسوب می‌شود. اگر شرایط محیطی جدیدی رخ دهد که در داده‌های آموزشی وجود نداشته باشد، مدل ممکن است رفتار نامناسب یا ناایمن از خود نشان دهد. این مسئله که به‌عنوان مشکل «انتقال توزیع%
\LTRfootnote{distribution shift}»% 
 شناخته می‌شود، به‌ویژه در محیط‌های دینامیکی و غیرقابل پیش‌بینی اهمیت بیشتری پیدا می‌کند.

محدودیت مهم دیگر یادگیری نظارتی، ماهیت غیرفعال آن است. در این چارچوب، مدل تنها از داده‌های از پیش جمع‌آوری‌شده یاد می‌گیرد و امکان تعامل فعال با محیط یا اصلاح رفتار بر اساس پیامد تصمیم‌ها را ندارد. به بیان دیگر، سیستم فاقد سازوکاری برای در نظر گرفتن هدف بلندمدت یا بهینه‌سازی رفتار در طول زمان است. مرورهای اخیر نیز تأکید می‌کنند که این ویژگی باعث می‌شود یادگیری نظارتی بیش‌تر برای تقلید رفتار یا پیش‌بینی محلی مناسب باشد تا تصمیم‌گیری تطبیقی در محیط‌های پویا و پیچیده%
\cite{alvarez2024DRLsurvey}.

در مجموع، یادگیری نظارتی نقش مهمی در استخراج الگوهای حرکتی و مدل‌سازی رفتارهای مرجع در ناوبری ربات‌ها ایفا می‌کند، اما به‌تنهایی پاسخ‌گوی نیازهای ناوبری خودران در محیط‌های دینامیکی و غیرقابل پیش‌بینی نیست. این محدودیت‌ها، انگیزه‌ای روشن برای بررسی روش‌هایی فراهم می‌کنند که بتوانند با تعامل مستقیم با محیط و در نظر گرفتن پیامد تصمیم‌ها در طول زمان، سیاست‌های ناوبری مناسب را یاد بگیرند. در بخش بعدی، نقش یادگیری بدون ناظر به‌عنوان یک رویکرد مکمل در ناوبری ربات‌های متحرک بررسی می‌شود.

\subsection{یادگیری بدون نظارت در ناوبری ربات‌های متحرک}

«یادگیری بدون نظارت%
\LTRfootnote{Unsupervised Learning}» %
به دسته‌ای از روش‌های یادگیری ماشین اطلاق می‌شود که در آن‌ها داده‌های آموزشی فاقد برچسب صریح بوده و هدف اصلی، کشف ساختارها و الگوهای نهفته در داده‌های خام است. در این چارچوب، برخلاف یادگیری نظارتی، خروجی مطلوب از پیش مشخص نشده و الگوریتم تلاش می‌کند با تحلیل توزیع داده‌ها، شباهت میان نمونه‌ها و روابط درونی آن‌ها، نمایش معناداری از داده‌های ورودی استخراج کند.

در حوزه ناوبری ربات‌های متحرک، یادگیری بدون نظارت به‌طور سنتی کمتر به‌عنوان یک راه‌حل مستقل برای تصمیم‌گیری حرکتی مورد استفاده قرار گرفته است. مرورهای انجام‌شده نشان می‌دهند که نقش غالب این روش‌ها بیش‌تر به‌عنوان یک ابزار کمکی در مراحل ادراک محیط، کاهش پیچیدگی داده‌های حسگری و یادگیری نمایش‌های فشرده از وضعیت ربات و محیط تعریف می‌شود. داده‌های حاصل از حسگرها معمولاً دارای ابعاد بالا، نویز قابل توجه و ساختار پیچیده هستند و استفاده مستقیم از آن‌ها در فرآیند تصمیم‌گیری می‌تواند منجر به افزایش هزینه محاسباتی و کاهش پایداری سامانه ناوبری شود. در چنین شرایطی، یادگیری بدون نظارت می‌تواند با استخراج ویژگی‌های معنادار، ورودی مناسبی برای مراحل بعدی تصمیم‌گیری فراهم کند%
\cite{Zhu2021,s25113394}. 

با این حال، در سال‌های اخیر پژوهش‌هایی منتشر شده‌اند که نشان می‌دهند یادگیری بدون نظارت می‌تواند نقشی فراتر از یک ابزار پیش‌پردازش ایفا کند. برای نمونه، در چارچوب %
 \lr{UP3}%
 ، یادگیری بدون نظارت به‌صورت مستقیم برای تولید مسیرهای حرکتی در مسئله ناوبری ربات‌های متحرک مورد استفاده قرار گرفته است. در این رویکرد، بدون استفاده از برچسب‌های دستی یا سیاست‌های مرجع، مدل تلاش می‌کند با بهره‌گیری از ساختار داده‌های مشاهده‌شده و پیش‌بینی تحول محیط، مسیرهای قابل‌اجرا و ایمن را در محیط‌های با قابلیت مشاهده محدود تولید کند. نتایج ارائه‌شده نشان می‌دهند که این روش قادر است در برخی سناریوهای پیچیده، مسیرهایی هموار و سازگار با محدودیت‌های حرکتی ربات تولید کند، بدون آن‌که نیازمند تعامل گسترده با محیط یا فرآیند آموزش مبتنی بر پاداش باشد%
\cite{Luo2025UP3}.

این دسته از پژوهش‌ها نشان می‌دهد که یادگیری بدون نظارت می‌تواند در شرایط خاص، به‌ویژه زمانی که دسترسی به داده‌های برچسب‌خورده یا طراحی تابع پاداش مناسب دشوار است، به‌عنوان یک گزینه جایگزین یا مکمل برای روش‌های دیگر مطرح شود. با این وجود، باید توجه داشت که چنین رویکردهایی معمولاً فرضیات مشخصی درباره ساختار محیط یا قابلیت پیش‌بینی آن دارند و تعمیم آن‌ها به محیط‌های کاملاً دینامیکی و غیرقابل پیش‌بینی همچنان با چالش مواجه است.

به‌طور کلی، یادگیری بدون نظارت در ناوبری ربات‌های متحرک بیش‌تر به‌عنوان رویکردی مکمل در کنار سایر روش‌های یادگیری و کنترل مورد استفاده قرار می‌گیرد. حتی در کارهایی که از این روش به‌صورت مستقیم برای تولید مسیر استفاده شده است، محدودیت‌هایی در زمینه تعمیم‌پذیری، تضمین ایمنی و سازگاری با تغییرات سریع محیط مشاهده می‌شود. این محدودیت‌ها موجب شده‌اند که تمرکز اصلی پژوهش‌های اخیر به سمت روش‌هایی معطوف شود که بتوانند با تعامل فعال با محیط و در نظر گرفتن پیامد تصمیم‌ها در طول زمان، سیاست‌های ناوبری مؤثرتر و تطبیق‌پذیرتری را یاد بگیرند. در بخش بعدی، یادگیری تقلیدی به‌عنوان رویکردی که تلاش می‌کند رفتار مرجع را مستقیماً در سیاست ناوبری وارد کند، مورد بررسی قرار می‌گیرد.

\subsection{یادگیری تقلیدی در ناوبری ربات‌های متحرک}

«یادگیری تقلیدی%
\LTRfootnote{Imitation Learning}»
به دسته‌ای از روش‌های یادگیری ماشین اطلاق می‌شود که در آن‌ها ربات تلاش می‌کند رفتار یک عامل مرجع، معمولاً انسان یا یک الگوریتم کنترلی از پیش طراحی‌شده، را با مشاهده و تقلید از نمونه‌های رفتاری بیاموزد. در این چارچوب، داده‌های آموزشی شامل دنباله‌ای از مشاهدات محیط و اعمال متناظر با آن‌ها هستند و هدف یادگیری، بازتولید رفتار مرجع در شرایط مشابه است. برخلاف یادگیری تقویتی، در یادگیری تقلیدی نیازی به تعریف تابع پاداش صریح وجود ندارد و رفتار مطلوب به‌طور ضمنی از طریق نمایش‌های ارائه‌شده مشخص می‌شود.

در حوزه ناوبری ربات‌های متحرک، یادگیری تقلیدی به‌عنوان رویکردی میانی میان یادگیری نظارتی و یادگیری تقویتی مطرح می‌شود. از یک سو، این روش مانند یادگیری نظارتی از داده‌های برچسب‌خورده استفاده می‌کند و از سوی دیگر، به‌طور مستقیم با مسئله تصمیم‌گیری و انتخاب عمل سروکار دارد. مرورهای کلاسیک در این حوزه نشان می‌دهند که یادگیری تقلیدی به‌ویژه برای یادگیری رفتارهای پیچیده‌ای که طراحی صریح آن‌ها دشوار است، مانند اجتناب از برخورد در محیط‌های شلوغ، گزینه‌ای جذاب محسوب می‌شود%
\cite{argall2009survey}.

یکی از رایج‌ترین پیاده‌سازی‌های یادگیری تقلیدی در ناوبری، تقلید رفتار یک سیاست مرجع است که معمولاً با استفاده از روش‌های کلاسیک برنامه‌ریزی یا کنترل تولید شده است. در این حالت، ربات با مشاهدهٔ دنباله‌های حرکتی تولیدشده توسط عامل مرجع، نگاشت میان وضعیت محیط و فرمان‌های حرکتی را می‌آموزد. این رویکرد امکان ترکیب دانش موجود در روش‌های کلاسیک با انعطاف‌پذیری مدل‌های یادگیری را فراهم می‌کند و می‌تواند منجر به تولید رفتارهایی هموار و قابل‌قبول در محیط‌های نسبتاً پیچیده شود.

در سال‌های اخیر، یادگیری تقلیدی به‌صورت انتها به انتها نیز مورد توجه قرار گرفته است. برای نمونه، در پژوهش «کادویلا%
 \LTRfootnote{Codevilla}»%
  و همکاران، یک چارچوب یادگیری تقلیدی شرطی ارائه شده است که در آن مدل با استفاده از داده‌های رانندگی انسانی، رفتار ناوبری را مستقیماً از داده‌های ادراکی می‌آموزد. اگرچه این کار در زمینه رانندگی خودکار انجام شده است، چارچوب آن به‌طور گسترده در ناوبری ربات‌های متحرک نیز مورد استفاده و ارجاع قرار گرفته و نشان می‌دهد که یادگیری تقلیدی می‌تواند رفتارهای پیچیده را بدون تعریف صریح قوانین کنترلی بازتولید کند%
\cite{codevilla2018endtoend}.

با وجود مزایای یادگیری تقلیدی، این روش با محدودیت‌های مهمی نیز همراه است. یکی از چالش‌های اساسی، وابستگی شدید عملکرد سیستم به کیفیت و تنوع داده‌های نمایشی است. اگر رفتار مرجع ناقص یا محدود به شرایط خاصی باشد، مدل یادگرفته‌شده نیز در مواجهه با شرایط جدید یا غیرمنتظره دچار افت عملکرد خواهد شد. این مسئله که به‌عنوان مشکل \lr{covariate shift} شناخته می‌شود، به‌ویژه در محیط‌های دینامیکی که وضعیت‌ها به‌سرعت تغییر می‌کنند، اهمیت بیشتری پیدا می‌کند.

محدودیت دیگر یادگیری تقلیدی، عدم وجود سازوکار صریح برای بهینه‌سازی رفتار بر اساس یک هدف بلندمدت است. از آنجا که مدل صرفاً رفتار مشاهده‌شده را تقلید می‌کند، امکان بهبود تدریجی عملکرد یا کشف راهبردهای جدید فراتر از داده‌های مرجع محدود است. مرورهای اخیر نیز نشان می‌دهند که اگرچه یادگیری تقلیدی می‌تواند نقطه شروع مناسبی برای یادگیری رفتارهای ناوبری باشد، اما به‌تنهایی برای دستیابی به رفتارهای تطبیق‌پذیر و بهینه در محیط‌های پویا کافی نیست%
\cite{Zhu2021}.


\section{یادگیری تقویتی به‌عنوان چارچوب تصمیم‌گیری در ناوبری ربات‌های متحرک}

یادگیری تقویتی به‌عنوان یکی از رویکردهای اصلی یادگیری ماشین، چارچوبی را برای یادگیری تصمیم‌گیری از طریق تعامل مستقیم با محیط فراهم می‌کند. در این رویکرد، به‌جای اتکا به داده‌های برچسب‌خورده یا تقلید صرف از رفتار مرجع، عامل یادگیرنده با انجام عمل در محیط و مشاهده پیامد آن، به‌تدریج رفتار خود را اصلاح می‌کند. این ویژگی، یادگیری تقویتی را به گزینه‌ای مناسب برای مسائلی تبدیل می‌کند که ماهیت آن‌ها ذاتاً تصمیم‌محور، وابسته به زمان و متأثر از پیامدهای بلندمدت هستند.

مسئله ناوبری ربات‌های متحرک نمونه‌ای شاخص از چنین مسائل تصمیم‌گیری است. ربات ناوبر در هر لحظه با انتخاب یک فرمان حرکتی، نه‌تنها وضعیت فعلی خود را تغییر می‌دهد، بلکه شرایط تصمیم‌گیری در گام‌های بعدی را نیز تحت تأثیر قرار می‌دهد. به بیان دیگر، تصمیم‌های حرکتی در ناوبری دارای وابستگی زمانی بوده و اثر آن‌ها محدود به یک گام زمانی کوتاه نیست. این ویژگی باعث می‌شود که رویکردهایی که تنها بر پیش‌بینی یا تقلید رفتارهای محلی تکیه دارند، در مواجهه با شرایط پیچیده و دینامیکی با محدودیت مواجه شوند.

در مقایسه با روش‌های مبتنی بر یادگیری نظارتی یا تقلیدی، یادگیری تقویتی امکان در نظر گرفتن هدف‌های بلندمدت را به‌صورت ضمنی فراهم می‌کند. در این چارچوب، رفتار مطلوب ربات از طریق تعریف معیارهای کلی عملکرد، مانند رسیدن به هدف، اجتناب از برخورد و حفظ همواری حرکت، شکل می‌گیرد و نیازی به تعیین صریح خروجی مطلوب برای هر وضعیت وجود ندارد. این موضوع به‌ویژه در محیط‌هایی که پیش‌بینی کامل رفتار موانع یا طراحی قوانین صریح دشوار است، اهمیت پیدا می‌کند%
\cite{Zhu2021}.

ویژگی مهم دیگر یادگیری تقویتی، ماهیت تعاملی آن است. عامل یادگیرنده با تجربه مستقیم محیط، قادر است راهبردهایی را کشف کند که لزوماً در داده‌های از پیش جمع‌آوری‌شده یا رفتارهای مرجع مشاهده نمی‌شوند. این توانایی کشف رفتارهای نوظهور، یادگیری تقویتی را به ابزاری مناسب برای مواجهه با شرایط غیرمنتظره و تغییرپذیر تبدیل کرده است؛ شرایطی که در بسیاری از سناریوهای ناوبری واقعی، به‌ویژه در محیط‌های دینامیکی و شلوغ، به‌طور طبیعی وجود دارند.

مرورهای اخیر نشان می‌دهند که به‌دلیل همین ویژگی‌ها، یادگیری تقویتی به‌تدریج به‌عنوان یکی از چارچوب‌های اصلی تصمیم‌گیری در ناوبری ربات‌های متحرک مطرح شده است. این رویکرد نه‌تنها امکان یادگیری سیاست‌های حرکتی تطبیق‌پذیر را فراهم می‌کند، بلکه زمینه‌ای را برای ترکیب اهداف مختلف ناوبری در قالب یک رفتار منسجم ایجاد می‌نماید. با این حال، استفاده مؤثر از یادگیری تقویتی در ناوبری مستلزم توجه به شرایط محیطی، قیود حرکتی و چالش‌های عملی است که در ادامه این بخش مورد بررسی قرار خواهند گرفت%
\cite{s25113394}.
\subsection{خانواده‌های الگوریتم‌های یادگیری تقویتی در ناوبری ربات‌های متحرک}

الگوریتم‌های یادگیری تقویتی مورد استفاده در ناوبری ربات‌های متحرک را می‌توان بر اساس نحوه‌ی نمایش سیاست و سازوکار یادگیری به چند خانواده‌ی اصلی تقسیم کرد. این دسته‌بندی در ادبیات پژوهشی به‌عنوان یک چارچوب مفهومی برای تحلیل مزایا و محدودیت‌های روش‌های مختلف  یادگیری تقویتی در مسائل ناوبری شناخته می‌شود و نقش مهمی در انتخاب الگوریتم مناسب برای سامانه‌های حرکتی واقعی ایفا می‌کند. مرورهای جامع اخیر نشان می‌دهند که هر یک از این خانواده‌ها، رفتار متفاوتی در مواجهه با پیچیدگی‌های محیط‌های دینامیکی از خود نشان می‌دهند%
\cite{Alsolami2025AutonomousNavigationSurvey,Le2024DRLNavigationSurvey}.
\begin{itemize}
\item{\textbf{روش‌های مبتنی بر ارزش:}}
نخستین خانواده، روش‌های «\textbf{مبتنی بر ارزش}%
\LTRfootnote{Value-based}» هستند که در آن‌ها سیاست به‌طور غیرمستقیم و از طریق یادگیری تابع ارزش یا تابع $Q$ استخراج می‌شود. الگوریتم‌هایی مانند 
\lr{Q-learning} و نسخه‌های عمیق‌تر آن نظیر DQN در این دسته قرار می‌گیرند. این روش‌ها عمدتاً برای مسائل با فضای عمل گسسته طراحی شده‌اند و در برخی سناریوهای ساده‌ی ناوبری نتایج قابل قبولی ارائه داده‌اند. با این حال، گسسته‌سازی فضای عمل در سامانه‌های حرکتی واقعی می‌تواند منجر به کاهش دقت کنترلی و بروز رفتارهای ناپیوسته شود؛ مسئله‌ای که کاربرد این خانواده را در پلتفرم‌های دارای دینامیک پیوسته محدود می‌کند%
\cite{Alsolami2025AutonomousNavigationSurvey}.

\item{\textbf{روش‌های مبتنی بر گرادیان سیاست:}}
خانواده‌ی دوم شامل روش‌های «\textbf{مبتنی بر گرادیان سیاست}%
\LTRfootnote{Policy gradient}» است که سیاست کنترلی را به‌صورت مستقیم و پارامتری یاد می‌گیرند. الگوریتم‌هایی نظیر REINFORCE، TRPO و PPO در این دسته قرار دارند. این روش‌ها امکان کار با فضای عمل پیوسته را فراهم می‌کنند و از این نظر برای مسائل ناوبری ربات‌های متحرک مناسب‌تر از روش‌های مبتنی بر ارزش هستند. با این حال، بسیاری از این الگوریتم‌ها از نوع on-policy بوده و برای یادگیری مؤثر نیازمند داده‌های جدید و همگن هستند؛ موضوعی که می‌تواند به کاهش کارایی نمونه و افزایش هزینه‌ی آموزش در محیط‌های پیچیده و دینامیکی منجر شود%
\cite{Kober2013RLRobotics,Le2024DRLNavigationSurvey}.

\item{\textbf{روش‌های کنشگر--منتقد:}}
سومین خانواده، روش‌های «\textbf{Actor--Critic}» هستند که ترکیبی از دو رویکرد قبلی را به‌کار می‌گیرند. در این چارچوب، بخش actor مسئول یادگیری سیاست و بخش critic مسئول برآورد کیفیت تصمیم‌ها از طریق تابع ارزش است. این ساختار امکان یادگیری پایدارتر در فضاهای حالت و عمل با ابعاد بالا را فراهم می‌کند و در سال‌های اخیر به‌طور گسترده در مسائل ناوبری ربات‌های متحرک مورد استفاده قرار گرفته است. الگوریتم‌های مختلفی در این خانواده معرفی شده‌اند که در پژوهش‌های جدید به‌عنوان گزینه‌های مناسب برای کنترل پیوسته و ناوبری در محیط‌های دینامیکی شناخته می‌شوند%
\cite{Alsolami2025AutonomousNavigationSurvey,Le2024DRLNavigationSurvey}.

یکی از ویژگی‌های مهم روش‌های ،actor--critic امکان استفاده از داده‌های off-policy و در نتیجه بهبود «\textbf{کارایی نمونه}» در مقایسه با بسیاری از روش‌های on-policy است. این ویژگی در مسائل ناوبری واقعی، که جمع‌آوری داده‌ی فیزیکی پرهزینه و همراه با ریسک است، اهمیت ویژه‌ای دارد. با این حال، این الگوریتم‌ها نیز با چالش‌هایی نظیر حساسیت به تنظیم هایپرپارامترها و پیچیدگی پیاده‌سازی همراه هستند که باید در طراحی سامانه‌های عملی مدنظر قرار گیرند.
\end{itemize}
به‌طور کلی، مرور ادبیات نشان می‌دهد که روند پژوهش‌ها در حوزه ناوبری ربات‌های متحرک، به‌تدریج از روش‌های مبتنی بر ارزش به سمت روش‌های actor--critic با قابلیت مدیریت فضای عمل پیوسته و قیود دینامیکی حرکت کرده است. این روند، زمینه‌ساز توسعه و به‌کارگیری چارچوب‌های پیشرفته‌ی یادگیری تقویتی در سامانه‌های حرکتی واقعی شده و مبنای تصمیم‌گیری در پژوهش‌های جدید را تشکیل می‌دهد%
\cite{Alsolami2025AutonomousNavigationSurvey}.
\subsection{الگوریتم‌های یادگیری تقویتی برای فضاهای عمل پیوسته در ناوبری}

در بسیاری از کاربردهای عملی ناوبری ربات‌های متحرک، فضای عمل به‌صورت ذاتی «\textbf{پیوسته}»%
\LTRfootnote{continuous action space}%
	است. برای مثال، فرمان‌های سرعت خطی و زاویه‌ای، گشتاور چرخ‌ها یا شتاب‌های اعمال‌شده به سامانه‌های حرکتی، همگی مقادیری پیوسته هستند. در چنین شرایطی، گسسته‌سازی فضای عمل، که در روش‌های مبتنی بر ارزش رایج است، می‌تواند منجر به کاهش دقت کنترلی، افزایش نوسانات حرکتی و رفتارهای غیرهموار شود. از این رو، بخش قابل توجهی از پژوهش‌های اخیر در ناوبری ربات‌ها به استفاده از الگوریتم‌های یادگیری تقویتی مناسب برای فضاهای عمل پیوسته معطوف شده است%
\cite{Kober2013RLRobotics,Alsolami2025AutonomousNavigationSurvey}.

یکی از نخستین تلاش‌های موفق در این زمینه، الگوریتم «\textbf{DDPG}%
\LTRfootnote{Deep Deterministic Policy Gradient}» است که به‌عنوان یک روش actor--critic با سیاست قطعی معرفی شد%
\cite{Lillicrap2016}. %
 این الگوریتم امکان یادگیری مستقیم سیاست‌های پیوسته را فراهم می‌کند و به‌طور گسترده در مسائل کنترل ربات‌های متحرک مورد استفاده قرار گرفته است. با این حال، مطالعات نشان داده‌اند که DDPG نسبت به نویز، تنظیم هایپرپارامترها و برآورد ناپایدار تابع ارزش حساس است، که این مسئله می‌تواند در محیط‌های دینامیکی ناوبری منجر به رفتارهای ناپایدار شود%
\cite{Kober2013RLRobotics}.

برای رفع بخشی از این محدودیت‌ها، الگوریتم «\textbf{\lr{TD3}}%
\LTRfootnote{Twin Delayed Deep Deterministic Policy Gradient}» معرفی شد که با استفاده از دو شبکه‌ی critic و به‌روزرسانی‌های تأخیری سیاست، مشکل بیش‌برآوردی تابع ارزش را کاهش می‌دهد%
\cite{FujimotoTD3}. %
 این بهبود باعث افزایش پایداری آموزش و کیفیت سیاست‌های یادگرفته‌شده در مقایسه با DDPG شده است و در برخی مطالعات ناوبری بدون نقشه عملکرد بهتری از خود نشان داده است. با این حال، سیاست‌های قطعی \lr{TD3} همچنان می‌توانند در محیط‌های بسیار پویا یا دارای عدم‌قطعیت بالا، از نظر اکتشاف فضای حالت با محدودیت مواجه شوند%
\cite{Nasti2025MaplessTD3}.

در مقابل، الگوریتم‌های مبتنی بر سیاست‌های «\textbf{تصادفی}%
\LTRfootnote{stochastic policies}» رویکرد متفاوتی را دنبال می‌کنند. در این میان، الگوریتم «\textbf{کنشگر-منتقد نرم}%
\LTRfootnote{Soft Actor--Critic-SAC}» با افزودن مؤلفه‌ی آنتروپی به تابع هدف، تلاش می‌کند میان بهره‌برداری از سیاست‌های آموخته‌شده و اکتشاف فضای حالت تعادل برقرار کند
\cite{Haarnoja2018SAC}. %
 این ویژگی به‌ویژه در محیط‌های ناوبری دینامیکی اهمیت دارد، زیرا ربات باید بتواند در مواجهه با موانع متحرک و تغییرات پیش‌بینی‌ناپذیر، رفتارهای متنوع و سازگار از خود نشان دهد. مطالعات مروری اخیر نشان می‌دهند که الگوریتم‌های مبتنی بر سیاست‌های تصادفی، از نظر پایداری آموزش و کارایی نمونه، مزایایی نسبت به سیاست‌های کاملاً قطعی دارند%
\cite{Le2024DRLNavigationSurvey,Alsolami2025AutonomousNavigationSurvey}.

به‌طور کلی، ادبیات پژوهشی نشان می‌دهد که حرکت از سیاست‌های گسسته به سیاست‌های پیوسته، و از روش‌های مبتنی بر ارزش به سمت چارچوب‌های actor--critic پیشرفته، یک روند غالب در ناوبری ربات‌های متحرک بوده است. این روند به‌ویژه در سامانه‌هایی با قیود دینامیکی، نیاز به حرکت هموار و تعامل ایمن با محیط، اهمیت بیشتری پیدا می‌کند و مبنای استفاده از الگوریتم‌های پیشرفته‌ی یادگیری تقویتی در پژوهش‌های اخیر را شکل می‌دهد%
\cite{Alsolami2025AutonomousNavigationSurvey}.

\subsection{مزایای یادگیری تقویتی نسبت به روش‌های پیشین ناوبری}

یکی از دلایل اصلی گرایش پژوهش‌های اخیر به یادگیری تقویتی در حوزه ناوبری، محدودیت‌های ذاتی روش‌های پیشین در مواجهه با محیط‌های پیچیده و پویا است. روش‌های کلاسیک ناوبری، هرچند از نظر محاسباتی کارآمد و قابل تحلیل هستند، اما معمولاً بر مدل‌سازی دقیق محیط یا فرضیات ساده‌کننده درباره رفتار موانع متکی‌اند. این وابستگی باعث می‌شود که عملکرد آن‌ها در شرایطی که محیط به‌طور مداوم تغییر می‌کند یا اطلاعات کامل در دسترس نیست، با افت مواجه شود. در مقابل، یادگیری تقویتی با تکیه بر تعامل مستقیم با محیط، چارچوبی انعطاف‌پذیرتر برای تصمیم‌گیری در چنین شرایطی فراهم می‌کند%
\cite{Wang2020G2RL}.

در مقایسه با روش‌های مبتنی بر یادگیری نظارتی، یادگیری تقویتی به داده‌های برچسب‌خورده وابسته نیست و رفتار مطلوب را از طریق بازخورد کلی عملکرد می‌آموزد. این ویژگی در مسائل ناوبری از اهمیت ویژه‌ای برخوردار است، زیرا تعیین خروجی بهینه برای تمام وضعیت‌های ممکن محیط عملاً غیرممکن است. در یادگیری تقویتی، به‌جای آموزش نگاشت مستقیم ورودی به خروجی، معیارهای کلی مانند ایمنی، کارایی و رسیدن به هدف تعریف می‌شوند و عامل یادگیرنده خود به کشف سیاست مناسب می‌پردازد. این رویکرد امکان یادگیری رفتارهایی را فراهم می‌کند که فراتر از داده‌های آموزشی اولیه هستند%
\cite{Zhu2021}.

یادگیری تقلیدی نیز اگرچه قادر است رفتارهای پیچیده را با استفاده از نمایش‌های مرجع بازتولید کند، اما معمولاً در تعمیم به شرایط جدید با محدودیت مواجه می‌شود. از آنجا که سیاست یادگرفته‌شده در این روش‌ها به شدت به داده‌های نمایشی وابسته است، مواجهه با وضعیت‌هایی که در داده‌های آموزشی وجود نداشته‌اند می‌تواند منجر به کاهش شدید عملکرد شود. یادگیری تقویتی با فراهم کردن امکان تعامل فعال با محیط و اصلاح تدریجی سیاست، توانایی سازگاری بیشتری در برابر تغییرات محیطی از خود نشان می‌دهد. این مزیت به‌ویژه در محیط‌های دینامیکی و شلوغ، که رفتار موانع قابل پیش‌بینی نیست، اهمیت بیشتری پیدا می‌کند%
\cite{s25113394}.

از منظر تصمیم‌گیری، یکی از مزایای مهم یادگیری تقویتی نسبت به بسیاری از روش‌های پیشین، توانایی در نظر گرفتن پیامدهای بلندمدت تصمیم‌ها است. در ناوبری ربات‌های متحرک، یک تصمیم حرکتی ممکن است در کوتاه‌مدت ایمن به نظر برسد، اما در گام‌های بعدی ربات را در موقعیتی نامطلوب قرار دهد. یادگیری تقویتی با هدف بیشینه‌سازی عملکرد تجمعی، امکان یادگیری راهبردهایی را فراهم می‌کند که تنها بر واکنش‌های محلی تکیه ندارند، بلکه توالی تصمیم‌ها را به‌صورت یکپارچه در نظر می‌گیرند. این ویژگی در مطالعات مختلف به‌عنوان یکی از برتری‌های اصلی این رویکرد نسبت به الگوریتم‌های واکنشی گزارش شده است%
\cite{Wang2020G2RL}.

علاوه بر این، یادگیری تقویتی قابلیت ادغام طبیعی با معیارهای متنوع عملکرد را فراهم می‌کند. در بسیاری از کاربردهای ناوبری، اهداف متعددی به‌طور هم‌زمان مطرح هستند؛ از جمله اجتناب از برخورد، کوتاه بودن مسیر، همواری حرکت و رعایت محدودیت‌های حرکتی ربات. چارچوب یادگیری تقویتی این امکان را می‌دهد که این معیارها به‌صورت یکپارچه در فرآیند یادگیری لحاظ شوند، بدون آن‌که نیاز به طراحی قوانین صریح و جداگانه برای هر یک وجود داشته باشد. مطالعات کاربردی نشان می‌دهند که این ویژگی می‌تواند منجر به تولید رفتارهای حرکتی منسجم‌تر و طبیعی‌تر نسبت به روش‌های مبتنی بر قواعد شود%
\cite{Lee2022DRLNavigation}.

در مجموع، مزایای یادگیری تقویتی نسبت به روش‌های کلاسیک، نظارتی و تقلیدی باعث شده است که این رویکرد به‌عنوان یکی از چارچوب‌های اصلی ناوبری ربات‌های متحرک مطرح شود. انعطاف‌پذیری در مواجهه با عدم قطعیت، توانایی یادگیری رفتارهای تطبیق‌پذیر و امکان در نظر گرفتن اهداف بلندمدت، از جمله عواملی هستند که یادگیری تقویتی را به گزینه‌ای جذاب برای مسائل ناوبری پیچیده تبدیل کرده‌اند. با این حال، همان‌گونه که در ادامه مورد بحث قرار خواهد گرفت، بهره‌گیری مؤثر از این مزایا مستلزم توجه به چالش‌ها و محدودیت‌های عملی این رویکرد است.
\subsection{چالش‌ها و محدودیت‌های یادگیری تقویتی در ناوبری}

یادگیری تقویتی عمیق علی‌رغم پیشرفت‌های چشمگیر در زمینه ناوبری ربات‌های متحرک، هنوز با چالش‌های مهمی در کاربردهای واقعی مواجه است. برخلاف روش‌های کلاسیک که رفتار سیستم به‌صورت صریح و قابل پیش‌بینی طراحی می‌شود، در DRL سیاست کنترلی از طریق تجربه و تعامل با محیط شکل می‌گیرد؛ این ویژگی اگرچه انعطاف‌پذیری بالایی ایجاد می‌کند، اما هم‌زمان منجر به بروز محدودیت‌هایی در آموزش و پیاده‌سازی عملی می‌شود%
\cite{Le2024DRLNavigationSurvey}.

یکی از چالش‌های اساسی، «کارایی نمونه%
\LTRfootnote{sample efficiency}»%
است. الگوریتم‌های DRL معمولاً برای همگرا شدن به سیاست‌های قابل قبول نیازمند حجم بسیار زیادی از داده‌های تعاملی هستند. این مسئله در شبیه‌سازی قابل مدیریت است، اما در محیط‌های واقعی می‌تواند منجر به افزایش هزینه، استهلاک سخت‌افزار و حتی بروز رفتارهای ناایمن در حین یادگیری شود. در محیط‌های دینامیکی که شامل موانع متحرک و تعامل با انسان هستند، این چالش اهمیت بیشتری پیدا می‌کند، زیرا امکان آزمون و خطای گسترده عملاً وجود ندارد%
\cite{Le2024DRLNavigationSurvey}.

چالش مهم دیگر پایداری و همگرایی الگوریتم‌ها در فضاهای حالت و عمل با ابعاد بالا است. ناوبری ربات‌های متحرک معمولاً با فضای حالت پیوسته، داده‌های حسگری پرنویز و محدودیت‌های دینامیکی همراه است. در چنین شرایطی، الگوریتم‌های DRL ممکن است دچار نوسانات شدید در فرآیند آموزش شوند یا به سیاست‌هایی همگرا شوند که تنها در شرایط خاص عملکرد مناسبی دارند. حساسیت بالا به تنظیم هایپرپارامترها و ساختار شبکه عصبی، این مشکل را تشدید می‌کند%
\cite{Le2024DRLNavigationSurvey}.

مسئلهٔ طراحی تابع پاداش یکی از چالش‌برانگیزترین جنبه‌های یادگیری تقویتی در ناوبری است. در مسائل واقعی، ربات باید به‌صورت هم‌زمان چندین هدف متضاد را رعایت کند؛ برای مثال، حرکت سریع‌تر ممکن است با افزایش خطر برخورد همراه باشد، در حالی که حرکت ایمن‌تر می‌تواند منجر به مسیرهای طولانی و ناکارآمد شود. ترکیب این اهداف در قالب یک تابع پاداش یکتا بدون ایجاد سوگیری یا رفتارهای ناخواسته، نیازمند طراحی دقیق و آزمون‌های گسترده است. در بسیاری از مطالعات گزارش شده است که تعریف نامناسب پاداش می‌تواند باعث یادگیری رفتارهای غیرایمن یا غیرکاربردی شود%
\cite{Le2024DRLNavigationSurvey}.

چالش قابل توجه دیگر انتقال از شبیه‌سازی به واقعیت است. اگرچه شبیه‌سازی ابزار قدرتمندی برای آموزش الگوریتم‌های DRL فراهم می‌کند، اما تفاوت‌های اجتناب‌ناپذیر میان شبیه‌سازی و دنیای واقعی، از جمله نویز حسگرها، مدل‌سازی ناقص دینامیک و رفتار غیرقابل پیش‌بینی موانع، می‌تواند باعث کاهش شدید عملکرد سیاست‌های آموزش‌دیده در محیط واقعی شود. این مسئله یکی از موانع اصلی در تجاری‌سازی و کاربرد صنعتی روش‌های DRL در ناوبری ربات‌ها محسوب می‌شود%
\cite{Le2024DRLNavigationSurvey}.

در سناریوهای «ناوبری بدون نقشه%
\LTRfootnote{mapless navigation}»، چالش‌های فوق تشدید می‌شوند. در این حالت، ربات بدون دسترسی به نقشهٔ از پیش‌ساخته‌شده باید تنها بر اساس داده‌های آنی حسگرها تصمیم‌گیری کند. این امر منجر به پاداش‌های پراکنده، دشواری در کاوش فضای حالت و افزایش احتمال همگرایی به سیاست‌های زیر‌بهینه می‌شود. مطالعات اخیر نشان داده‌اند که استفاده از الگوریتم‌های actor–critic بهبود‌یافته، مانند نسخه‌های اصلاح‌شدهٔ \lr{TD3} می‌تواند تا حدی پایداری یادگیری و کیفیت سیاست‌ها را در چنین شرایطی افزایش دهد، هرچند چالش‌های بنیادین همچنان باقی می‌مانند%
\cite{Nasti2025MaplessTD3}.

در مجموع، چالش‌هایی مانند کارایی نمونه، پایداری آموزش، طراحی تابع پاداش و انتقال به محیط‌های واقعی نشان می‌دهند که یادگیری تقویتی عمیق، علی‌رغم توانمندی‌های بالقوه، هنوز نیازمند توسعه و تکامل بیشتر برای کاربرد مطمئن در ناوبری ربات‌های متحرک در محیط‌های دینامیکی است. این چالش‌ها زمینه‌ساز پژوهش‌های گسترده‌ای در سال‌های اخیر بوده‌اند و انگیزهٔ اصلی برای توسعه چارچوب‌های ترکیبی و الگوریتم‌های پیشرفته‌تر محسوب می‌شوند.
\subsection{رویکردهای ایمن و ترکیبی در ناوبری مبتنی بر یادگیری تقویتی}

با وجود توانمندی‌های یادگیری تقویتی عمیق در تولید سیاست‌های کنترلی تطبیقی، یکی از مهم‌ترین دغدغه‌های مطرح در ادبیات ناوبری ربات‌های متحرک، مسئله‌ی \textbf{ایمنی} و رعایت قیود فیزیکی در طول فرآیند یادگیری و اجرا است. برخلاف روش‌های کلاسیک کنترل که قیود سیستم به‌صورت صریح در طراحی لحاظ می‌شوند، الگوریتم‌های یادگیری تقویتی ذاتاً فاقد تضمین‌های ایمنی هستند و در صورت عدم اعمال ملاحظات اضافی، ممکن است به رفتارهای ناایمن یا غیرقابل پیش‌بینی منجر شوند%
\cite{Le2024DRLNavigationSurvey}.

\begin{itemize}
\item{\textbf{یادگیری تقویتی ایمن و مقید:}}
یکی از رویکردهای مطرح در ادبیات، «\textbf{یادگیری تقویتی ایمن}%
\LTRfootnote{Safe Reinforcement Learning}» یا «\textbf{یادگیری تقویتی مقید}%
\LTRfootnote{Constrained Reinforcement Learning}» است که در آن، قیود ایمنی به‌صورت صریح در فرآیند یادگیری وارد می‌شوند. این قیود می‌توانند شامل محدودیت‌هایی بر فاصله از موانع، سرعت و شتاب مجاز، یا سایر پارامترهای مرتبط با ایمنی حرکت باشند. در این چارچوب‌ها، هدف تنها بیشینه‌سازی پاداش نیست، بلکه رعایت قیود نیز به‌عنوان بخشی از مسئله‌ی تصمیم‌گیری در نظر گرفته می‌شود%
\cite{Alsolami2025AutonomousNavigationSurvey}.

\item{\textbf{طراحی پاداش با در نظر گرفتن ایمنی و کیفیت حرکت:}}
رویکرد رایج‌تر در بسیاری از کاربردهای ناوبری، استفاده از «\textbf{شکل‌دهی پاداش}%
\LTRfootnote{reward shaping}» برای اعمال ملاحظات ایمنی و کیفی است. در این روش، به‌جای تعریف قیود صریح، مؤلفه‌های تنبیهی در تابع پاداش گنجانده می‌شوند تا رفتارهایی مانند نزدیک شدن بیش از حد به موانع، تغییرات ناگهانی فرمان یا نوسانات شدید حرکتی کاهش یابند. این نوع طراحی پاداش به‌ویژه برای سامانه‌های حرکتی واقعی که نرمی حرکت، راحتی و پایداری اهمیت دارند، مورد توجه قرار گرفته است%
\cite{Le2024DRLNavigationSurvey}.

\item{\textbf{روش‌های ترکیبی مبتنی بر یادگیری تقویتی:}}
دسته‌ی دیگری از پژوهش‌ها به سمت «\textbf{روش‌های ترکیبی}%
\LTRfootnote{hybrid approaches}» حرکت کرده‌اند که در آن‌ها یادگیری تقویتی با روش‌های کلاسیک برنامه‌ریزی مسیر یا کنترل ترکیب می‌شود. در این رویکردها، یادگیری تقویتی معمولاً نقش تصمیم‌گیری سطح بالا یا تولید سیاست کلی را بر عهده دارد، در حالی که تضمین ایمنی، رعایت قیود دینامیکی و اجرای دقیق فرمان‌ها به ماژول‌های کلاسیک مانند برنامه‌ریزها یا کنترل‌کننده‌های پیش‌بین سپرده می‌شود. این ترکیب امکان بهره‌گیری از انعطاف‌پذیری RL و در عین حال حفظ پایداری و ایمنی روش‌های کلاسیک را فراهم می‌کند%
\cite{Alsolami2025AutonomousNavigationSurvey}.
\end{itemize}
به‌طور کلی، مرور ادبیات نشان می‌دهد که برای کاربردهای واقعی ناوبری، به‌ویژه در پلتفرم‌هایی با جرم، ابعاد و محدودیت‌های دینامیکی قابل توجه، استفاده از یادگیری تقویتی به‌صورت خام کمتر مورد استقبال قرار گرفته و اغلب با ملاحظات ایمنی، طراحی پاداش هدفمند یا چارچوب‌های ترکیبی همراه بوده است. این رویکردها بیانگر تلاش پژوهشگران برای کاهش فاصله‌ی میان قابلیت‌های بالقوه‌ی یادگیری تقویتی و الزامات عملی سامانه‌های حرکتی واقعی هستند.
\subsection{مقایسه مطالعات منتخب در حوزه ناوبری ربات‌های متحرک}
به‌منظور جمع‌بندی منسجم مطالعات مرورشده، جدول \ref{tab:nav_literature_comparison} مقایسه‌ای میان مهم‌ترین خانواده‌های روش‌های کلاسیک، روش‌های مبتنی بر یادگیری و مطالعات نزدیک به مسئله پژوهش حاضر ارائه می‌دهد. هدف از این جدول، صرفاً فهرست کردن منابع نیست، بلکه روشن کردن نقش هر گروه از روش‌ها، نقاط قوت آن‌ها و محدودیت‌های آن‌ها نسبت به مسئله ناوبری اسکوتر چهار چرخ محرک در محیط دینامیکی است.

\begingroup
\scriptsize
\setlength{\tabcolsep}{3pt}
\renewcommand{\arraystretch}{1.35}
\newcolumntype{P}[1]{>{\raggedright\arraybackslash}p{#1}}
\newcolumntype{C}[1]{>{\centering\arraybackslash}p{#1}}
\newcolumntype{R}[1]{>{\raggedleft\arraybackslash}p{#1}}

\begin{longtable}{|C{2.35cm}|P{2.35cm}|P{2.65cm}|P{3.45cm}|P{3.95cm}|}
	\caption{مقایسه رویکردها و مطالعات منتخب در حوزه ناوبری ربات‌های متحرک}
	\label{tab:nav_literature_comparison}\\
	\hline
	\centering\textbf{مرجع} &
	\centering\textbf{رویکرد یا الگوریتم} &
	\centering\textbf{نوع مسئله یا محیط} &
	\centering\textbf{نقطه قوت و نقش در ادبیات} &
	\centering\arraybackslash\textbf{محدودیت نسبت به پژوهش حاضر} \\
	\hline
	\endfirsthead
	
	\hline
	\textbf{مرجع} &
	\textbf{رویکرد یا الگوریتم} &
	\textbf{نوع مسئله یا محیط} &
	\textbf{نقطه قوت و نقش در ادبیات} &
	\textbf{محدودیت نسبت به پژوهش حاضر} \\
	\hline
	\endhead
	
	\hline
	\endfoot
	
	\cite{Dijkstra1959,Hart1968AStar,koenig2005dstar} &
	روش‌های جست‌وجومحور مانند \lr{Dijkstra}، \lr{A*} و \lr{D* Lite} &
	برنامه‌ریزی مسیر سراسری در محیط‌های دارای نقشه &
	تولید مسیر کم‌هزینه در نمایش گرافی محیط و مناسب برای محیط‌های ایستا یا نیمه‌ایستا &
	مسیر تولیدشده عمدتاً هندسی است و به‌تنهایی قیود دینامیکی، محدودیت ترمز، موانع متحرک و تصمیم‌گیری لحظه‌ای را پوشش نمی‌دهد. \\
	\hline
	
	\cite{Kavraki1996PRM,LaValle1998RRT,Karaman2011RRTstar} &
	روش‌های نمونه‌برداری مانند \lr{PRM}، \lr{RRT} و \lr{RRT*} &
	برنامه‌ریزی در فضای پیکربندی، به‌ویژه در محیط‌های پیچیده &
	توانایی یافتن مسیر در فضاهای پیکربندی پیچیده بدون نیاز به گسسته‌سازی کامل محیط &
	مسیرها معمولاً نیازمند هموارسازی و پردازش تکمیلی هستند و برای تصمیم‌گیری بلادرنگ در محیط‌های شدیداً دینامیکی کافی نیستند. \\
	\hline
	
	\cite{khatib1986,Koren1991PotentialFields} &
	میدان پتانسیل مصنوعی \lr{APF} &
	اجتناب محلی از مانع و حرکت به سمت هدف &
	سادگی محاسباتی، پیاده‌سازی آسان و قابلیت اجرای بلادرنگ &
	دارای مشکل کمینه محلی، حساس به چینش موانع و فاقد لحاظ صریح قیود دینامیکی و رفتار زمانی موانع متحرک است. \\
	\hline
	
	\cite{Fiorini1998VO,VanDenBerg2008RVO,VanDenBerg2011ORCA} &
	روش‌های مبتنی بر سرعت مانند \lr{VO}، \lr{RVO} و \lr{ORCA} &
	اجتناب از برخورد در حضور موانع متحرک و عامل‌های چندگانه &
	تحلیل برخورد در فضای سرعت و توجه به زمان برخورد، نه فقط فاصله لحظه‌ای &
	وابسته به فرض حرکت کوتاه‌مدت قابل پیش‌بینی، تخمین دقیق سرعت موانع و معمولاً محدود به اجتناب محلی از برخورد است. \\
	\hline
	
	\cite{fox1997dwa} &
	روش پنجره دینامیکی \lr{DWA} &
	برنامه‌ریزی محلی در فضای سرعت &
	لحاظ سرعت، شتاب و توان توقف در انتخاب فرمان و امکان اجرای بلادرنگ &
	افق تصمیم‌گیری کوتاه دارد و عملکرد آن به تنظیم تابع هزینه وابسته است؛ در محیط‌های شلوغ ممکن است رفتار محافظه‌کارانه یا نوسانی ایجاد کند. \\
	\hline
	
	\cite{Quinlan1993ElasticBands,Rosmann2013TEB,Rosmann2015TEB} &
	نوار کشسان زمان‌بندی‌شده \lr{TEB} &
	برنامه‌ریزی محلی مبتنی بر بهینه‌سازی &
	ورود زمان‌بندی حرکت، محدودیت سرعت، شتاب و فاصله از موانع در بهینه‌سازی مسیر محلی &
	به وزن‌دهی تابع هزینه، کیفیت نقشه محلی، داده‌های حسگری و توان بازبهینه‌سازی سریع در محیط‌های دینامیکی وابسته است. \\
	\hline
	
	\cite{rawlings2009mpc,Bemporad1999,Mayne2005TubeMPC} &
	کنترل پیش‌بینانه مبتنی بر مدل \lr{MPC} و \lr{Tube MPC} &
	کنترل و برنامه‌ریزی مقید در سامانه‌های دینامیکی &
	توانایی لحاظ هم‌زمان قیود حرکتی، کنترلی، ایمنی و اهداف بهینه‌سازی در افق پیش‌بینی &
	نیازمند مدل نسبتاً دقیق، پیش‌بینی مناسب محیط و حل مسئله بهینه‌سازی در هر گام زمانی است؛ در محیط‌های شلوغ و نامطمئن هزینه طراحی و محاسبات بالا می‌رود. \\
	\hline
	
	\cite{s23249672} &
	مقایسه روش‌های سنتی و مبتنی بر یادگیری تقویتی عمیق &
	ناوبری داخلی در سناریوهای دینامیکی &
	نشان دادن تفاوت عملکرد روش‌های کلاسیک و \lr{DRL} در شرایط دینامیکی و وابستگی انتخاب روش به ویژگی محیط &
	بیشتر نقش مقایسه‌ای دارد و چارچوب اختصاصی برای پلتفرم‌هایی با قیود حرکتی جدی مانند اسکوتر ارائه نمی‌کند. \\
	\hline
	
	\cite{pfeiffer2018nav} &
	برنامه‌ریزی انتهابه‌انتها مبتنی بر داده &
	ناوبری ربات زمینی از ادراک تا تصمیم &
	کاهش وابستگی به طراحی دستی ماژول‌های جداگانه و اتصال مستقیم ادراک به تصمیم حرکتی &
	وابستگی بالا به داده آموزشی و تعمیم‌پذیری محدود در سناریوهای خارج از توزیع؛ قیود دینامیکی پلتفرم به‌صورت مرکزی بررسی نمی‌شود. \\
	\hline
	
	\cite{argall2009survey,codevilla2018endtoend} &
	یادگیری تقلیدی و راندن شرطی انتهابه‌انتها &
	یادگیری سیاست حرکتی از داده یا نمایش خبره &
	مناسب برای استفاده از تجربه انسانی یا سیاست مرجع و کاهش نیاز به طراحی صریح قوانین تصمیم‌گیری &
	کیفیت سیاست به داده خبره وابسته است و در مواجهه با شرایط دیده‌نشده ممکن است دچار افت عملکرد شود؛ بهینه‌سازی بلندمدت پاداش مانند \lr{RL} در مرکز روش نیست. \\
	\hline
	
	\cite{Luo2025UP3} &
	برنامه‌ریزی پیش‌بین بدون‌نظارت \lr{UP3} &
	محیط‌های ناقصاً مشاهده‌پذیر &
	توجه به پیش‌بینی آینده و تصمیم‌گیری در شرایطی که مشاهده محیط کامل نیست &
	مستقیماً چارچوب کنترل پیوسته مبتنی بر پاداش برای پلتفرم دارای قیود دینامیکی ارائه نمی‌کند. \\
	\hline
	
	\cite{Wang2020G2RL} &
	یادگیری تقویتی با راهنمایی سراسری &
	برنامه‌ریزی مسیر در محیط‌های دینامیکی &
	ترکیب اطلاعات برنامه‌ریزی سراسری با تصمیم‌گیری یادگیری‌محور برای بهبود ناوبری &
	هرچند برای محیط‌های دینامیکی مفید است، تمرکز اصلی آن بر پلتفرم‌هایی مانند اسکوتر چهار چرخ با قیود نرمی حرکت، توقف و راحتی سرنشین نیست. \\
	\hline
	
	\cite{Nasti2025MaplessTD3} &
	ناوبری بدون نقشه با الگوریتم بهبودیافته \lr{TD3} &
	ناوبری \lr{Mapless} در فضای عمل پیوسته &
	نزدیک به مسئله پژوهش حاضر از نظر استفاده از یادگیری تقویتی عمیق و فضای عمل پیوسته &
	تمرکز اصلی بر ناوبری بدون نقشه است و ابعاد خاص پلتفرم اسکوتر، ایمنی سرنشین، نرمی فرمان و قیود حرکتی سنگین‌تر به‌صورت محوری بررسی نمی‌شود. \\
	\hline
	
	\cite{Pacini2024RLwheelchair,Cui2025OutdoorWheelchair} &
	ناوبری و کمک‌رانش برای پلتفرم‌های کمک‌حرکتی مانند ویلچر هوشمند &
	پلتفرم‌های انسانی/کمک‌حرکتی در محیط واقعی یا نیمه‌واقعی &
	از نظر ماهیت پلتفرم و اهمیت ایمنی کاربر به مسئله اسکوتر خودران نزدیک‌تر از ربات‌های کوچک آزمایشگاهی است &
	معمولاً بر معماری کمک‌رانندگی، نقشه هزینه یا اجتناب از مانع تمرکز دارد و چارچوب کامل یادگیری سیاست پیوسته برای اسکوتر چهار چرخ محرک را پوشش نمی‌دهد. \\
	\hline
	
\end{longtable}

\endgroup

بررسی جدول \ref{tab:nav_literature_comparison} نشان می‌دهد که روش‌های کلاسیک ناوبری هر یک بخشی از مسئله را به‌خوبی پوشش می‌دهند. روش‌های جست‌وجومحور و نمونه‌برداری برای تولید مسیر سراسری مفیدند، روش‌هایی مانند \lr{VO}، \lr{ORCA} و \lr{DWA} برای اجتناب محلی از برخورد کاربرد دارند، و روش‌های بهینه‌سازی‌محور مانند \lr{TEB} و \lr{MPC} امکان لحاظ برخی قیود حرکتی و زمانی را فراهم می‌کنند. با این حال، این روش‌ها معمولاً به مدل‌سازی صریح محیط، تنظیم تابع هزینه، کیفیت نقشه محلی، پیش‌بینی حرکت موانع یا مدل دینامیکی وابسته‌اند و در محیط‌های دینامیکی، شلوغ و ناقصاً مشاهده‌پذیر به‌تنهایی پاسخ کاملی ارائه نمی‌کنند.

از سوی دیگر، مطالعات مبتنی بر یادگیری نشان می‌دهند که روش‌های داده‌محور و یادگیری تقویتی عمیق می‌توانند بخشی از محدودیت‌های روش‌های کلاسیک را کاهش دهند؛ زیرا سیاست تصمیم‌گیری از طریق داده یا تعامل با محیط شکل می‌گیرد و امکان واکنش به سناریوهای پیچیده‌تر فراهم می‌شود. با وجود این، بخش قابل توجهی از این مطالعات یا بر ربات‌های متحرک ساده‌تر تمرکز دارند، یا در محیط‌های شبیه‌سازی‌شده با فرضیات ساده‌شده ارزیابی شده‌اند، یا قیود خاص پلتفرم‌هایی مانند اسکوتر چهار چرخ محرک، از جمله اینرسی، محدودیت توقف، نرمی فرمان و ایمنی سرنشین را به‌صورت مرکزی بررسی نکرده‌اند.

بنابراین، شکاف اصلی ادبیات موجود در کمبود چارچوبی دیده می‌شود که ناوبری خودران را برای پلتفرمی مانند اسکوتر چهار چرخ محرک، در محیط دینامیکی و با فضای عمل پیوسته، به‌صورت هم‌زمان از منظر رسیدن به هدف، اجتناب از برخورد، رعایت قیود حرکتی، تولید فرمان‌های نرم و حفظ ایمنی بررسی کند. این شکاف، ضرورت استفاده از رویکردهای یادگیری تقویتی عمیق و به‌ویژه الگوریتم‌های مناسب برای کنترل پیوسته را در پژوهش حاضر توجیه می‌کند.

\section{جمع‌بندی مرور ادبیات و تبیین شکاف پژوهشی}

مرور انجام‌شده در این فصل و مقایسه مطالعات منتخب در جدول \ref{tab:nav_literature_comparison} نشان داد که مسئله ناوبری خودران ربات‌های متحرک، به‌ویژه در محیط‌های دینامیکی، همچنان یکی از موضوعات فعال و چندوجهی در حوزه رباتیک و هوش مصنوعی است. روش‌های کلاسیک ناوبری، علی‌رغم بلوغ نظری، قابلیت تحلیل‌پذیری و پیاده‌سازی نسبتاً روشن، معمولاً به مدل‌سازی صریح محیط، تنظیم تابع هزینه، کیفیت نقشه محلی یا پیش‌بینی حرکت موانع وابسته‌اند. این وابستگی باعث می‌شود که عملکرد آن‌ها در محیط‌های شلوغ، ناقصاً مشاهده‌پذیر و دارای موانع متحرک، به‌ویژه برای پلتفرم‌هایی با قیود حرکتی جدی، با محدودیت مواجه شود.

در مقابل، روش‌های مبتنی بر یادگیری، به‌ویژه یادگیری تقویتی عمیق، امکان یادگیری سیاست‌های تصمیم‌گیری غیرخطی و تطبیقی را از طریق تعامل با محیط فراهم می‌کنند. این ویژگی باعث شده است که در سال‌های اخیر از این روش‌ها برای ناوبری بدون نقشه، اجتناب از برخورد، تصمیم‌گیری در محیط‌های دینامیکی و کنترل در فضای عمل پیوسته استفاده شود. با این حال، بررسی ادبیات نشان می‌دهد که کاربرد یادگیری تقویتی در سامانه‌های واقعی ناوبری همچنان با چالش‌هایی مانند کارایی نمونه، پایداری آموزش، طراحی تابع پاداش، ایمنی رفتار یادگرفته‌شده و تعمیم‌پذیری به شرایط دیده‌نشده همراه است.

نکته مهم‌تر آن است که بخش قابل توجهی از مطالعات موجود بر ربات‌های متحرک سبک، محیط‌های ساده‌سازی‌شده، سناریوهای شبیه‌سازی‌شده یا پلتفرم‌هایی با حرکت محدود تمرکز داشته‌اند. در چنین مطالعاتی، معیارهایی مانند رسیدن به هدف و اجتناب از برخورد معمولاً در مرکز ارزیابی قرار می‌گیرند، در حالی که ملاحظاتی مانند نرمی فرمان، محدودیت شتاب و ترمز، اینرسی پلتفرم، پایداری حرکت و راحتی یا ایمنی سرنشین کمتر به‌صورت هم‌زمان و مسئله‌محور بررسی می‌شوند. این موضوع برای پلتفرمی مانند اسکوتر چهار چرخ محرک اهمیت بیشتری دارد، زیرا تصمیم ناوبری در چنین سامانه‌ای صرفاً به انتخاب یک مسیر هندسی یا جلوگیری از برخورد محدود نمی‌شود، بلکه باید به تولید فرمان‌های پیوسته، نرم، ایمن و قابل اجرا نیز منجر شود.

از این منظر، شکاف اصلی ادبیات موجود را می‌توان در کمبود چارچوب‌هایی دانست که ناوبری خودران را برای پلتفرم‌هایی با دینامیک و قیود اجرایی جدی، مانند اسکوتر چهار چرخ محرک، در محیط‌های دینامیکی و با فضای عمل پیوسته بررسی کنند. چنین چارچوبی باید بتواند میان چند هدف هم‌زمان، از جمله رسیدن به هدف، اجتناب از برخورد، رعایت محدودیت‌های حرکتی، کاهش تغییرات ناگهانی فرمان و حفظ ایمنی حرکت، تعادل برقرار کند. بنابراین، مسئله پژوهش حاضر نه صرفاً یافتن مسیر، بلکه طراحی و ارزیابی یک سیاست تصمیم‌گیری پیوسته برای ناوبری ایمن و هموار اسکوتر خودران در محیط دینامیکی است.

بر این اساس، شکاف‌های پژوهشی اصلی که این پایان‌نامه بر آن‌ها متمرکز است، به‌صورت زیر خلاصه می‌شوند:
\begin{itemize}
	\item کمبود مطالعات مسئله‌محور درباره ناوبری خودران پلتفرم‌هایی مانند اسکوتر چهار چرخ محرک که دارای اینرسی، محدودیت توقف و قیود حرکتی جدی هستند؛
	\item تمرکز محدود بسیاری از پژوهش‌های موجود بر ارزیابی هم‌زمان موفقیت در رسیدن به هدف، اجتناب از برخورد، نرمی حرکت و قابلیت اجرای فرمان‌ها؛
	\item وابستگی بسیاری از روش‌های کلاسیک به مدل‌سازی صریح محیط، تنظیم تابع هزینه یا پیش‌بینی حرکت موانع در محیط‌های دینامیکی؛
	\item چالش باقی‌مانده در طراحی سیاست‌های یادگیری تقویتی برای فضای عمل پیوسته که بتوانند میان کارایی ناوبری، ایمنی و نرمی فرمان تعادل برقرار کنند.
\end{itemize}

با توجه به این شکاف‌ها، توسعه و ارزیابی یک چارچوب ناوبری خودران مبتنی بر یادگیری تقویتی برای اسکوتر چهار چرخ محرک، در محیط دینامیکی و تحت قیود عملی حرکت، ضرورت پیدا می‌کند. این چارچوب باید بتواند علاوه بر دستیابی به هدف و اجتناب از برخورد، کیفیت فرمان‌های تولیدشده و سازگاری آن‌ها با محدودیت‌های حرکتی پلتفرم را نیز در نظر بگیرد. بر همین اساس، در فصل بعدی، فرمول‌بندی مسئله، مدل‌سازی محیط، تعریف حالت و عمل، طراحی پاداش و روش پیشنهادی پژوهش تشریح می‌شود.